\[
\textbf{Theorem.}\quad \sqrt{5}\text{ is irrational.}
\]

\[
\textbf{Proof.}
\]
Suppose, for contradiction, that \(\sqrt{5}\) is rational. By definition of rationality, there exist integers \(m,n\) with \(n\neq 0\) such that

\[
\sqrt{5}=\frac{m}{n}.
\]

We now justify that we may choose such a representation in lowest terms. Let

\[
d=\gcd(m,n).
\]

Since \(n\neq 0\), \(d\ge 1\). Because \(d\mid m\) and \(d\mid n\), there exist integers \(a,b\) such that

\[
m=da,\qquad n=db.
\]

Then \(b\neq 0\), and

\[
\frac{m}{n}=\frac{da}{db}=\frac{a}{b}.
\]

Moreover, by dividing \(m\) and \(n\) by their greatest common divisor, we obtain

\[
\gcd(a,b)=1.
\]

Thus we may write

\[
\sqrt{5}=\frac{a}{b},
\]

where \(a,b\in \mathbb{Z}\), \(b\neq 0\), and \(\gcd(a,b)=1\).

Squaring both sides gives

\[
5=\frac{a^2}{b^2}.
\]

Multiplying by \(b^2\), which is nonzero, gives

\[
a^2=5b^2.
\]

Hence

\[
5\mid a^2.
\]

We claim that \(5\mid a\). To see this without appealing to any unstated theorem, consider the possible remainders of \(a\) modulo \(5\). Every integer is congruent modulo \(5\) to one of

\[
0,1,2,3,4.
\]

Squaring these residues gives

\[
0^2\equiv 0,\quad 1^2\equiv 1,\quad 2^2\equiv 4,\quad 3^2\equiv 9\equiv 4,\quad 4^2\equiv 16\equiv 1 \pmod 5.
\]

Thus the only way \(a^2\equiv 0\pmod 5\) is if

\[
a\equiv 0\pmod 5.
\]

Therefore \(5\mid a\). Hence there exists an integer \(k\) such that

\[
a=5k.
\]

Substituting this into \(a^2=5b^2\), we get

\[
(5k)^2=5b^2.
\]

So

\[
25k^2=5b^2.
\]

Dividing both sides by \(5\) gives

\[
5k^2=b^2.
\]

Therefore

\[
5\mid b^2.
\]

Using the same residue argument as above, \(5\mid b^2\) implies \(5\mid b\).

Thus \(5\mid a\) and \(5\mid b\). This contradicts the assumption that

\[
\gcd(a,b)=1,
\]

because if both \(a\) and \(b\) are divisible by \(5\), then their greatest common divisor is at least \(5\).

Therefore our supposition that \(\sqrt{5}\) is rational must be false. Hence

\[
\sqrt{5}\text{ is irrational.}
\]

\[
\boxed{\sqrt{5}\notin \mathbb{Q}}
\]